\begin{table*}[p]
\scriptsize
\centering
\caption{\textbf{Analytic model families used in the main analyses.} The table summarizes descriptive and conversation-level model families. Adjusted models are interpreted as covariate-adjusted associations in naturalistic data, not as causal treatment-effect estimates.}
\label{tab:model_specs}
\setlength{\tabcolsep}{4pt}
\renewcommand{\arraystretch}{1.12}
\begin{tabularx}{\textwidth}{>{\raggedright\arraybackslash}p{0.19\textwidth}>{\raggedright\arraybackslash}p{0.20\textwidth}>{\raggedright\arraybackslash}p{0.28\textwidth}>{\raggedright\arraybackslash}X}
\toprule
Analysis family & Unit and outcome & Main contrast or model & Scope \\
\midrule
Engagement prevalence & User turns and conversations; cognitive, constructive and emotional engagement indicators. & Descriptive proportions by platform-by-task setting, task ecology and learning intent. & Establishes where learning-oriented engagement is visible; no causal contrast is implied. \\
Engagement composition & Cognitively engaged user turns; passive, active and constructive depth levels. & Composition of cognitive engagement within each platform-by-task setting. & Motivates constructive engagement as the focal strict outcome while preserving the broader cognitive-engagement hierarchy. \\
Interaction depth & Conversations; at least one constructive user turn. & Probability of constructive engagement across behavioural-depth buckets. & Depth is interpreted as interactional opportunity and task development, not as a manipulated exposure. \\
Scaffolding presence & Conversations; constructive-turn count and any constructive turn. & Poisson generalized linear models for counts and logistic models for binary outcomes, comparing conversations with versus without scaffolded assistant turns. & Adjusted estimates reduce measured compositional differences but remain associational because support is selected into by task difficulty and user state. \\
\bottomrule
\end{tabularx}
\end{table*}

\begin{table*}[p]
\scriptsize
\centering
\caption{\textbf{Additional analytic model families and robustness checks.} These analyses qualify the main associations by support form, temporal scale and available WildChat model metadata.}
\label{tab:model_specs_additional}
\setlength{\tabcolsep}{4pt}
\renewcommand{\arraystretch}{1.12}
\begin{tabularx}{\textwidth}{>{\raggedright\arraybackslash}p{0.19\textwidth}>{\raggedright\arraybackslash}p{0.20\textwidth}>{\raggedright\arraybackslash}p{0.28\textwidth}>{\raggedright\arraybackslash}X}
\toprule
Analysis family & Unit and outcome & Main contrast or model & Scope \\
\midrule
Support-form associations & Scaffolded conversations; cognitive or constructive engagement ratios. & Mean differences for conversations containing a given support-means label versus scaffolded conversations without that label, stratified by learning intent. & Descriptive because means labels are non-exclusive and can co-occur in the same assistant turn. \\
Local temporal coupling & Adjacent assistant--user or user--assistant turn pairs; next constructive user turn or next scaffolded assistant turn. & Conditional probability contrasts around adjacent turns; detailed estimands are listed in Supplementary Table~\ref{tab:temporal_estimands}. & Interpreted as local ordering and state-dependent coupling, not as randomized support effects. \\
Coarse temporal contrasts & Conversations containing scaffolded support; constructive engagement before and after the first scaffolded turn. & Mean after-minus-before contrast and first-half versus second-half robustness summaries where available. & Tests longer-span accumulation patterns, but the first scaffolded turn may be selected by difficulty, uncertainty or an already developing exchange. \\
WildChat model robustness & WildChat conversations or turns; main prevalence, association and temporal outcomes. & Repetition across user-facing model snapshots or coarse model families. & Robustness summaries only; model labels are observational and may be confounded with time, user mix and task mix. \\
\bottomrule
\end{tabularx}
\end{table*}
